%% IEEE conference paper -- compile with pdfLaTeX + BibTeX.
%% Double-blind submission. Uses refs.bib.
\documentclass[conference]{IEEEtran}
\IEEEoverridecommandlockouts

\usepackage{cite}
\usepackage{amsmath,amssymb,amsfonts}
\usepackage{graphicx}
\usepackage{booktabs}
\usepackage{balance}
\usepackage{xcolor}
\usepackage{url}
\usepackage{algorithm}
\usepackage{tikz}
\usetikzlibrary{positioning, shapes.geometric, arrows.meta, fit, calc}

\newcommand{\y}{\checkmark}
\newcommand{\n}{$\times$}
\newcommand{\figorbox}[3][\columnwidth]{%
  \IfFileExists{#2}{\includegraphics[width=#1]{#2}}%
  {\fbox{\parbox[c][3.2cm][c]{0.92\columnwidth}{\centering \texttt{#3}}}}}

\begin{document}

\title{How Much Real PHY Does a 5G Digital Twin Need?\\Two Constants, Measured on OpenAirInterface}

\author{\IEEEauthorblockN{}%
\IEEEauthorblockA{}}

\maketitle
\bstctlcite{BSTcontrol}

\begin{abstract}
Network digital twins are where O-RAN rApps are validated before they touch production, so the twin's link-to-system (L2S) abstraction---the map from SINR to throughput---decides which controllers look good. Received wisdom says this map must be measured on a real physical layer. We ask \emph{how much} of a real PHY a twin needs, and answer by measuring two constants on OpenAirInterface (OAI). Our open, CPU-only, multi-cell twin (OpenStreetMap $\to$ Sionna~RT $\to$ per-RE SINR $\to$ EESM $\to$ OAI's \texttt{nr\_dlsim} curve) hosts a traffic-steering rApp and three bridges to a live OAI 5G Standalone stack. We separate a twin's \emph{reporting} error from its \emph{decision} error, the latter measured as rApp \emph{regret}: the utility lost by planning on a surrogate twin and living in the OAI-grounded one. An uncalibrated Shannon curve inflates throughput $39.7\%$ and costs $0.767$ nats; a \emph{fixed} $2$\,dB margin halves the first and barely touches the second; a single \emph{fitted} margin $\delta{=}4.71$\,dB collapses both, to $1.6\%$ and $0.033$ nats. The customary ``decisions differ'' statistic ($25$--$38\%$ for every surrogate) separates none of them; only regret does. On the live stack the twin's link rate exceeds what OAI's MAC schedules by a \emph{constant} factor ($\kappa{=}0.470\pm0.004$, CV $0.8\%$ over MCS 0--28), so it cancels in ratios: the twin's \emph{relative} per-UE steering-gain prediction matches the real MAC to $1.2\pm5.7$ points ($95\%$ CI straddling zero). Delivered goodput lies a further $80\%$ down and grows with SINR---the emulator's ceiling, not the twin's error. Finally, a null control on two \emph{identical} real cells shows the natural goodput probe reporting $-21.0\pm1.4\%$ where truth is zero, retracting a multi-cell gap we previously reported.
\end{abstract}

\begin{IEEEkeywords}
Digital twin, 5G NR, OpenAirInterface, link-to-system abstraction, O-RAN, rApp, traffic steering.
\end{IEEEkeywords}

%=============================================================================
\section{Introduction}
%=============================================================================
The digital twin is becoming the environment in which AI/ML-driven RAN control---O-RAN rApps and xApps---is trained and validated before it touches production~\cite{dt_survey}. Its value hinges on data fidelity: a traffic-steering rApp decides which cell serves each user from signal, interference and load, which depend jointly on 3D propagation, on neighbouring cells, and on a link model turning SINR into a real modulation-and-coding outcome. Evaluating such a controller needs all three, across many layouts: our own single-seed pilot suggested a uniformly large rApp gain, which 24 runs tempered into a robust-direction, variable-magnitude result.

Open tooling splits into camps that do not meet: system-level simulators scale to many cells but abstract the PHY with idealized SINR-to-throughput curves, while OAI-based emulators execute the real 5G-NR stack but, with realistic propagation, cover a single link. No open platform combines all four of a real protocol stack, a site-specific ray-traced channel, a multi-cell network, and a closed control loop.

A subtler gap concerns the \emph{L2S abstraction} itself. System-level evaluation cannot afford full PHY processing per user, so it maps SINR to a modulation-and-coding outcome. Open simulators use idealized or 3GPP-calibrated curves; we obtain ours from OAI's real receiver (\texttt{nr\_dlsim}, LDPC decoding and rate matching). At 10\,dB SINR it yields 2.16\,bits/s/Hz against Shannon's 3.46---a 38\% implementation loss---with non-uniform MCS steps spanning up to 2.2\,dB.

The obvious claim is that this curve is therefore necessary. We tried to make it, and the evidence did not support it. The useful question is not \emph{whether} an idealized curve is wrong---it plainly is---but \emph{which part} of the real curve a controller is sensitive to. That needs a metric this literature does not use: whether two twins pick different actions is not whether picking the wrong one costs anything. We measure \emph{regret}, the utility an rApp loses by planning on one twin and being evaluated on another. Under regret, almost all of the OAI curve's value collapses onto one implementation margin $\delta$. A second constant, a derate $\kappa$, converts the twin's link rate into what a real MAC schedules. The twin needs these two numbers, and for this controller nothing else.

This paper contributes:
\begin{enumerate}
\item \emph{Two constants, and evidence that they are enough.} We separate a twin's \emph{reporting} error (what it says) from its \emph{decision} error (what acting on it costs, measured as rApp regret). An uncalibrated Shannon curve inflates throughput $39.7\%$ and costs $0.767$ nats of proportional-fair utility; a \emph{fixed} $2$\,dB margin halves the first and barely moves the second; a single \emph{fitted} margin $\delta{=}4.71$\,dB collapses both ($1.6\%$, $0.033$ nats). It is $\delta$'s value, not the presence of a margin, that a controller is sensitive to, and the curve's non-uniform shape buys nothing further. The customary ``fraction of layouts whose decisions differ'' is $25$--$38\%$ for \emph{every} surrogate and cannot distinguish them; only regret can.
\item \emph{A constant, emulator-independent MAC derate.} On a live OAI Standalone stack at matched bandwidth (106 PRB, 30\,kHz), the twin's link-level rate overpredicts what OAI's MAC schedules by $\kappa^{-1}$ with $\kappa{=}0.470\pm0.004$---constant to $0.8\%$ across MCS 0--28, of which the TDD duty $0.743$ is exactly computable. Constancy is the point: $\kappa$ cancels in ratios, so the twin's \emph{relative} per-UE steering-gain prediction matches the real MAC to $1.2\pm5.7$ points, while delivered goodput sits a further $80\%$ down and grows with SINR---the emulator's compute ceiling, not the twin's error. Conflating the two is what makes ``fidelity gap'' numbers irreproducible.
\item An open, CPU-only, multi-cell/multi-user digital twin and three twin-to-real-stack bridges, the third injecting \emph{per-cell} ray-traced channels on a two-DU stack---exploiting the UE's role as rfsimulator server, so the DUs' waveforms are physically summed: combined-waveform inter-cell interference, not a noise proxy.
\item \emph{A null control, and a retraction.} Give two real cells the \emph{same} channel and the true effect is exactly zero; the natural goodput probe then reports $-21.0\pm1.4\%$. We retract a multi-cell gap we previously reported, and state plainly that a controlled multi-cell measurement does not close on this emulator.
\item A 24-run, 3-city evaluation with bootstrap CIs, paired significance tests reported per scene as well as pooled, and an EESM-$\beta$ robustness control. We also retract a correlation: neither baseline outage nor peak cell load predicts rApp gain once computed \emph{within} scene, where both reverse sign.
\end{enumerate}

%=============================================================================
\section{Related Work}
%=============================================================================
\textbf{Open RAN digital-twin platforms.} No open platform combines a real protocol stack, a ray-traced channel, a multi-cell network and a closed control loop. Tiny-Twin~\cite{tinytwin} runs a real OAI stack with a RIC loop, single-cell from replayed traces. RIC-TaaP/ns-O-RAN and OpenTwin~\cite{opentwin} close rApp loops across many cells---OpenTwin adds twin-to-real KPM self-correction---but on ns-3's abstracted PHY. OAI+FlexRIC~\cite{oaiho} closes handover loops over the air, without a ray-traced channel. Commercial twins close the loop but are closed-source.

\textbf{Positioning against OWDT.} OWDT~\cite{owdt} is the closest work: it also pairs OAI with Sionna~RT and a real PHY. The distinction is architectural. (i)~OWDT is \emph{single-cell}: with one transmitter there is no inter-cell interference term, so its per-RE SINR reduces to an SNR and the association problem---which cell serves which user---does not exist. Our KPI engine carries a load-scaled interference sum over all non-serving cells (Eq.~\ref{eq:sinr}), and our third bridge makes that interference physically real by summing two DUs' waveforms rather than approximating it as a noise-floor rise. (ii)~OWDT is \emph{monitor-only}: it streams KPIs over E2 but nothing acts on them, so it cannot ask what a controller would have decided; a closed loop needs multi-cell KPIs \emph{and} a controller consuming them. (iii)~Consequently OWDT can report that a real PHY changes a KPI, but not what that change costs a decision---the question its design cannot pose.

\textbf{Link-to-system abstraction.} System-level simulation~\cite{tr36814} relies on L2S mapping: a link-level simulator characterises SINR-to-BLER per MCS, and an effective-SINR mapping (EESM~\cite{eesm,eesm_orig}) collapses a frequency-selective SINR vector into an AWGN-equivalent value. Our lookup is measured from OAI's own PHY rather than from idealised curves. Our interference-as-noise-floor bridge is a site-specific generalisation of the 3GPP OFDMA channel noise generator (OCNG)~\cite{ocng}: we derive each UE's interference from the ray-traced neighbour channels.

\textbf{O-RAN control loops.} The O-RAN architecture~\cite{oran,orantraffic} externalises RAN control: xApps on the near-RT RIC (via E2) and rApps on the non-RT RIC (via O1/A1). Traffic steering via cell-individual offsets (CIO) is a canonical rApp use case~\cite{dt_survey}, which a digital twin lets one validate before deployment.


%=============================================================================
\section{Digital Twin Architecture}
%=============================================================================
The platform is a pipeline with two consumers of one physics substrate: a fast analytical layer covering every cell and UE, and a real-stack layer grounding sampled links in OAI. Both are driven by the same ray-traced CFR: OSM geometry $\to$ Mitsuba scene $\to$ Sionna~RT $\to$ CFR $\to$ \{KPI engine $\to$ rApp\} and \{tap extraction $\to$ OAI rfsimulator\}. A WGS84 bounding box drives an OpenStreetMap query; footprints are extruded into a 3D Mitsuba scene with ITU materials. Cells and UEs come from a swappable source---a seeded random generator ($N$ sites $\times$ 3 sectors, uniform UE drops), or a CSV reader for real cell-site coordinates---so the twin runs on real sites without code changes. Sionna~RT~\cite{sionnart} places one transmitter per cell (3GPP sector pattern, planar array, downtilt) and one receiver per UE, and computes the channel frequency response of every cell-to-UE link over an OFDM resource grid; propagation settings and their rationale are in the reproducibility statement.

%=============================================================================
\section{OAI-Grounded KPI Engine}
\label{sec:kpi}
%=============================================================================
Let $\mathbf{H}_{u,c}(k)\in\mathbb{C}^{N_r\times N_t}$ be the ray-traced channel of link $c\!\to\!u$ on resource element (RE) $k$, $P^{\mathrm{tx}}_c$ the cell transmit power, $P^{\mathrm{rx}}_{u,c}=P^{\mathrm{tx}}_c\,\overline{\lVert\mathbf{H}_{u,c}\rVert_F^2}$ the wideband received power. Each UE associates to the cell maximising it, biased by a cell-individual offset (CIO)---the O-RAN control knob:
\begin{equation}
s(u)=\arg\max_{c}\; P^{\mathrm{rx}}_{u,c}\cdot 10^{\,\mathrm{CIO}_c/10}.
\label{eq:assoc}
\end{equation}
Each cell serves one UE at a time under airtime sharing, so the optimal single-stream precoder is the matched filter and the serving per-RE gain is $\sigma_{\max}$. With flat per-RE transmit power $p_c$, neighbour load $\rho\in[0,1]$, thermal noise $N_0=kTB\cdot\mathrm{NF}$:
\begin{equation}
\gamma_u(k)=\frac{p_{s}\,\sigma_{\max}\!\big(\mathbf{H}_{u,s}(k)\big)^2}
{\rho\sum_{c\neq s}\frac{p_c}{N_t}\lVert\mathbf{H}_{u,c}(k)\rVert_F^2+N_0}.
\label{eq:sinr}
\end{equation}
The per-RE SINRs are compressed to one effective SINR by exponential effective SINR mapping (EESM), preserving the ray tracer's frequency selectivity:
\begin{equation}
\gamma^{\mathrm{eff}}_u=-\beta\,\ln\!\Big(\tfrac{1}{K}\textstyle\sum_{k}
e^{-\gamma_u(k)/\beta}\Big),
\label{eq:eesm}
\end{equation}
with $\beta$ set per modulation order~\cite{eesm_orig} and one scale parameter. These are literature defaults, not values fitted to \texttt{nr\_dlsim} under frequency-selective channels; rather than assume that is benign we bound it in Sec.~\ref{sec:regret}. The bound is on magnitude, not on structural error across modulation orders.

\textbf{OAI-measured link curve.} The effective SINR maps to an MCS and spectral efficiency, $(m_u,\eta_u)=\Phi_{\mathrm{OAI}}(\gamma^{\mathrm{eff}}_u)$, through a curve \emph{measured offline from OAI's \texttt{nr\_dlsim}} (real LDPC decoding and rate matching), so it carries a real receiver's implementation losses: decoding margins, rate-matching overhead, non-uniform SINR spacing between MCS transitions.

At 10\,dB the OAI curve delivers 2.16\,bits/s/Hz---38\% below Shannon's 3.46---and the QPSK$\to$16QAM boundary spans 2.2\,dB against ${\sim}1$\,dB for adjacent QPSK steps. Whether that \emph{shape} matters to a controller, as opposed to the $38\%$ level, is what Sec.~\ref{sec:regret} answers.


On a static full-buffer snapshot, proportional-fair scheduling reduces to equal airtime, so the per-UE throughput is
\begin{equation}
R_u=\frac{B_{s}}{K_{s}}\,\eta_u\,(1-\mathrm{BLER}_{\mathrm{tgt}}),
\label{eq:rate}
\end{equation}
where $B_s$ is the serving cell bandwidth and $K_s$ its number of attached UEs.

%=============================================================================
\section{Traffic-Steering rApp}
%=============================================================================
The rApp adjusts the CIO vector to maximise the proportional-fair utility $U(\mathrm{CIO})=\sum_u\log R_u$. Because $U$ is piecewise-constant in the CIOs (it changes only when a UE's serving cell flips), a fixed small step stalls on plateaus. We therefore do coordinate ascent: starting from $\mathrm{CIO}=\mathbf{0}$, sweep the cells, line-searching each cell's CIO over a grid while every \emph{other} cell is held at its current best value (a true coordinate ascent, not a reset-to-zero sweep), and keep the maximiser of $U$; stop when no cell's line search improves $U$. This is monotonically non-decreasing in $U$ by construction. Raising a cell's CIO pulls boundary UEs in (absorbing load); lowering it sheds them to neighbours.

A sweep costs $C(|G|{-}1)$ utility evaluations ($C{=}9$, $G{=}\{-12..{+}12\}$\,dB at $1$\,dB, $|G|{=}25$): $216$ per sweep, $433$ over the $2$--$3$ sweeps to convergence, $4.7$\,s per layout ($10.7$\,ms each). The rApp is a \emph{reference} controller, not a contribution: it exposes what the L2S abstraction costs a decision. Any controller consuming per-UE rates would serve.


Every KPI evaluation reuses the same engine with a \texttt{cio\_db} vector, so the control loop rides on the exact twin used for the baseline. Its output is a set of per-UE handover decisions a real stack would execute via O1/telnet; it converges in 2--3 sweeps, suitable for the non-RT RIC timescale.

%=============================================================================
\section{Twin-to-Real-Stack Bridges}
\label{sec:bridges}
%=============================================================================
Three bridges connect the analytical twin to a live OAI stack. Their status, stated up front: Bridge~1 is validated and carries every measured fidelity number here; Bridge~2 is a design specification, \emph{not} a demonstrated result; Bridge~3's injection is verified on the stack, but the multi-cell measurement it enables does not close (Sec.~\ref{sec:multignb}). All measured fidelity results reflect the single-gNB link dimension only.

\textbf{Bridge 1: channel injection.} A link's ray-traced CFR becomes a sample-spaced impulse response (IFFT); the strongest causal taps are kept, unit-energy normalised, and mapped to sample indices. They override the rfsimulator's channel through a patch to OAI (\texttt{oai\_rt\_inject\_channel} in \texttt{random\_channel.c}), confirmed in the gNB log; the link then runs its real PHY/MAC over the site-specific channel.

\textbf{Bridge 2: interference-as-noise-floor.} A single-gNB link has no notion of the other cells a UE hears. Per UE, the twin's inter-cell interference becomes a noise-floor rise $\Delta N_u=10\log_{10}((I_u+N_0)/N_0)$, $I_u$ the load-scaled sum of non-serving cells' received power, written as a per-client OAI noise power---a site-specific generalisation of OCNG~\cite{ocng}.

\textbf{Per-cell channel injection.} With two DUs the noise-floor proxy is unnecessary. In the handover topology the UE is the rfsimulator \emph{server}, so the $k$-th DU to connect is filtered by channel model \texttt{rfsimu\_channel\_ue}$k$ and \texttt{rxAddInput} accumulates every connection into the UE's receive buffer. One taps block per DU therefore presents each cell's own ray-traced link gain, and the non-serving cell's waveform is summed in as genuine co-channel interference. We also make DU1 co-channel with DU0 (identical but for PCI/cell identity); the stock configuration places them on different carriers.

%=============================================================================
\section{Experimental Setup}
%=============================================================================
\textbf{Scenes.} Berlin-Mitte (2{,}623 buildings), Paris-Op\'era (121), Manhattan-Midtown (119); geometry fetched/built once per scene and reused across seeds.

\textbf{Random layouts.} Eight seeded layouts per scene---24 runs ($323$\,s on one CPU server). Radio and topology: 3.5\,GHz, $B=38.16$\,MHz (106 PRB at 30\,kHz SCS, matching the live OAI cell), 46\,dBm per cell, $4\times1$ BS array and single UE antenna, 7\,dB receiver noise figure, ray-tracing interaction depth 3, 9 cells and 30 UEs per layout, default inter-cell load $\rho=1.0$. The CFR is exported over all 1272 subcarriers, so the per-RE SINR and its EESM compression cover the cell's full band.

\textbf{Configuration sweep.} Neighbour load factor in $\{0,0.25,0.5,0.75,1.0\}$ on a cached channel per scene.

\textbf{Metrics.} Per-UE downlink throughput (median, mean); \emph{served cell-edge} rate, the 5th percentile among UEs with non-zero throughput (a raw 5th percentile collapses to zero once outage exceeds $5\%$); \emph{outage}, the fraction of UEs with zero throughput; Jain's fairness index; and \emph{peak cell load reduction}, the decrease in the maximum number of UEs attached to any single cell (negative $=$ the rApp sheds load, in UEs). Gains are reported with seeded bootstrap $95\%$ CIs and paired tests; a gain is undefined when its baseline is zero.

\textbf{Clustering.} The eight seeds of a scene share one geometry, so the pooled $n{=}24$ tests have a cluster structure and overstate the effective sample size. We therefore report the three per-scene tests ($n{=}8$ independent layouts each) alongside the pooled one, and treat per-scene directional consistency, not the pooled $p$ alone, as the evidence.


%=============================================================================
\section{Results}
\label{sec:results}
%=============================================================================
\subsection{Traffic-steering rApp}
Across 24 runs, relative to strongest-cell association (Table~\ref{tab:results}): the rApp helps the users a strongest-cell policy under-serves and sheds peak load. The median per-UE rate rises from $15.8$ to $18.5$\,Mbps and the served cell-edge rate from $5.8$ to $7.4$\,Mbps; both are significant under a paired Wilcoxon signed-rank test on the pooled absolute rates ($p{=}2.0{\times}10^{-4}$, $4.0{\times}10^{-4}$), as is the reduction in peak cell load ($p{=}5.7{\times}10^{-3}$). Because seeds within a scene share geometry, we also test per scene ($n{=}8$): median rate reaches $p{<}0.05$ in Manhattan and Paris ($0.016$) but not Berlin ($0.148$); served cell-edge in Berlin and Manhattan ($0.047$, $0.031$) but not Paris ($0.156$). The improvement is directionally consistent across all three independent geometries; the pooled $p$ corroborates, it is not $24$ independent samples. Three honest observations. Aggregate rate rises only $3.7\%$; Jain's index ($+6.2\%$) misses significance ($p{=}0.052$) and is \emph{negative} in Paris---the proportional-fair objective trades sum rate for the tail, and does not always win. Outage is bit-for-bit identical in all 24 runs, so no test applies: steering relocates covered users, it does not create coverage. And on $2$--$5$ layouts the rApp moves no UE, an exactly zero gain; those ties are excluded from the sign test (Dixon--Mood), not counted as failures.

Per scene: Berlin, the densest, has the lowest outage ($10.4\%$) and sheds the most peak load ($-1.0$ UEs) but gains least in median rate; Paris has the highest outage ($25.8\%$) and loses fairness; Manhattan gains most in fairness ($+17.5\%$). The direction is robust across all three for median and cell-edge rate, the magnitude strongly deployment-dependent---exactly what a single-scenario study would misrepresent.

\begin{table}[t]
\caption{rApp gains vs.\ strongest-cell association, 24 runs. Mean with a seeded percentile-bootstrap $95\%$ CI; $p$ from a paired Wilcoxon signed-rank test on the absolute rates (sign test for Jain); ``$+$/$-$/t'' = positive/negative/tied runs.}
\label{tab:results}
\centering
\footnotesize
\setlength{\tabcolsep}{4pt}
\begin{tabular}{@{}lccc@{}}
\toprule
Metric & Overall ($95\%$ CI) & $p$ & $+$/$-$/t \\
\midrule
Median tput & $+10.9$ $[6.8, 14.8]$ & $2{\cdot}10^{-4}$ & 19/3/2 \\
Served edge & $+25.9$ $[15.8, 36.9]$ & $4{\cdot}10^{-4}$ & 16/3/5 \\
Jain index  & $+6.2$ $[1.0, 11.6]$  & $0.052$ & 16/6/2 \\
Peak load (UE) & $-0.62$ $[-0.96, -0.29]$ & $6{\cdot}10^{-3}$ & --- \\
Mean tput & $+3.7$ $[1.2, 6.4]$ & $0.036$ & 15/7/2 \\
Outage & $0.0$ (all runs) & n/a & 0/0/24 \\
\bottomrule
\end{tabular}
\end{table}

\subsection{What a twin says vs.\ what acting on it costs}
\label{sec:regret}
The 38\% link-curve gap of Sec.~\ref{sec:kpi} is real, but does a controller care? Hold geometry, ray tracing, association and interference fixed; change \emph{only} the link curve. The \emph{reporting} error is how far a surrogate twin's predicted throughput sits from the OAI-grounded twin's. The \emph{decision} error is \emph{regret}: take the CIO the rApp selects on the surrogate, evaluate it on the OAI-grounded twin, compare with the CIO it would have chosen had it planned there: $\mathcal{R} = U(\mathrm{CIO}^{\star}_{\Phi_{\mathrm{OAI}}} \mid \Phi_{\mathrm{OAI}}) - U(\mathrm{CIO}^{\star}_{\Phi_{\mathrm{sur}}} \mid \Phi_{\mathrm{OAI}}) \ge 0$,
with $U$ the proportional-fair objective the rApp maximises. Reporting error is what you get wrong by reading a number off the twin; regret is what you pay for acting on it. We compare four surrogates over the same MCS/SE set, differing only in the SINR at which each MCS becomes available: the Shannon-optimal staircase; a \emph{fixed} $2$\,dB implementation margin, as a link-budget abstraction would assume without link-level data; a \emph{fitted} SINR offset $\Phi(\gamma)=\log_2(1+\gamma/\delta)$, $\delta=4.71$\,dB; and a fitted attenuation $\Phi(\gamma)=a\log_2(1+\gamma)$, $a=0.606$. The last three are what a practitioner calibrates with one parameter, and are the honest baselines---the Shannon staircase is a straw man no simulator uses.

The two errors turn out to be close to independent (Fig.~\ref{fig:regret}). The idealized curve inflates throughput by $39.7\%$ $[30.4, 49.9]$ and costs $0.767$ nats. A \emph{fixed} $2$\,dB margin halves the reporting error ($+24.4\%$) but barely touches the regret ($0.732$): assuming a margin is nearly worthless. The \emph{fitted} margin $\delta$ removes both---$+1.6\%$ $[-0.1, 3.5]$ and $0.033$ $[0.006, 0.071]$ nats, a $23\times$ reduction---so it is the \emph{value} of $\delta$, not a margin's presence, that a controller is sensitive to. The attenuation model reports well ($-5.7\%$) yet is as costly to act on as no calibration at all ($0.800$ nats): reporting accuracy does not imply decision quality. Most tellingly, the statistic this literature usually reports---the fraction of layouts on which the two twins select different steering actions---is $37.5\%$, $29.2\%$, $25.0\%$ and $37.5\%$ across the four surrogates. It cannot tell the good one from the bad ones. Nor can the set-overlap (Jaccard) statistic it is usually paired with, which has no natural null---its value depends on how many UEs happen to move. Regret has one: zero. What the twin needs from a real PHY is therefore $\delta$; the curve's non-uniform MCS spacing, which is what measuring OAI buys beyond one scalar, does not change what the controller does.

As a control on the twin's own free parameter, perturbing the EESM $\beta$ scale by $\pm20\%$ shifts the mean effective SINR by only $\mp0.15$\,dB, changes the rApp's decisions on $4$--$8\%$ of layouts, and costs at most $0.046$ nats---the same order as the residual regret of the best-calibrated curve, and ${\sim}17\times$ below an uncalibrated one. The conclusions above are not an artefact of that knob.

\begin{figure}[t]
\centering
\figorbox[0.68\columnwidth]{fig_decision_regret.pdf}{fig\_decision\_regret.pdf}
\caption{Reporting error (left) and the cost of acting on it (right), per surrogate twin, over 24 runs. A single fitted implementation margin $\delta$ collapses both; a fitted attenuation collapses only the first.}
\label{fig:regret}
\end{figure}

\subsection{The real fidelity gap}
\label{sec:realgap}
The analysis above compares two \emph{models}; we now close the loop on the real stack. A full OAI 5G Standalone system runs on the same CPU host---core (AMF/SMF/UPF), gNB, and a UE that registers, establishes a PDU session and carries user-plane traffic. On a single gNB we ground the \emph{link} dimension: pin the downlink operating point by capping the scheduler MCS over a clean channel, sweep the cap across MCS~0--28, and measure end-to-end goodput at each fixed MCS, which indexes the same OAI curve the twin uses (a shared SINR axis). Two details decide the answer. The live cell is 106\,PRB at 30\,kHz SCS, i.e.\ $B=38.16$\,MHz occupied, not the twin scenario's $20$\,MHz, so the twin's rate must be evaluated at the cell's own bandwidth. And a saturating UDP probe is unusable: in reverse mode \texttt{iperf3} reports the \emph{sending} rate, so with the scheduler pinned to MCS~4 (a $9.8$\,Mbps MAC) an 8\,Mbps offer reads exactly $8.00$\,Mbps at ${\sim}0\%$ loss. TCP cannot exceed the channel but needs tens of seconds to leave slow start, so we settle and measure its plateau.

Two references matter, and conflating them is what makes ``fidelity gap'' numbers irreproducible. Against \textbf{what OAI's MAC schedules}---real 5G-NR scheduler code paying TDD duty, DMRS, PDCCH, HARQ---the twin's link rate $B\,\eta_{\mathrm{OAI}}(1{-}\mathrm{BLER})$ is high by a \emph{constant} factor: $\kappa = 0.470\pm0.004$ ($52.9\%$ inflation, ranging only $52.4$--$53.6\%$ over MCS 0--28; CV $0.8\%$). Against \textbf{delivered application goodput} it is high by a median $80\%$, and that gap is \emph{not} constant: $72\%$ at MCS~0 to $87\%$ at MCS~28 (CV $27\%$). At MCS~28 the twin predicts $190.8$\,Mbps, TDD duty allows $141.7$, OAI's MAC schedules $88.9$, the application receives $24.9$. Only the first two steps are modelling: the TDD duty ($104/140 = 0.743$, deterministic) and the MAC's physical-channel and scheduler overhead ($0.633$; DMRS, PDCCH, HARQ) account for the \emph{whole} twin-to-MAC discrepancy---which is why $\kappa$ is constant. Everything below the MAC line is an \emph{upper bound} on unmodelled stack effects, not an explained quantity: it mixes the rfsimulator's compute ceiling, same-host contention and transport, which one machine cannot separate. The obvious probe for the transport term is unavailable---in reverse mode \texttt{iperf3}'s UDP report is the \emph{sending} rate, so through a $9.8$\,Mbps MAC an $8$\,Mbps offer reads $8.00$\,Mbps at ${\sim}0\%$ loss. We quote $\kappa$, not the $80\%$ (Fig.~\ref{fig:realgap}, left).

Constancy is not a detail: it \emph{is} the second constant. A multiplicative $\kappa$ cancels in any ratio, so it should not perturb the twin's \emph{relative} predictions. We test that by propagating each curve through the rApp over eight layouts (240 per-UE evaluations, 19 steered UEs), indexing every curve by the twin's own operating MCS---the twin's rate is a staircase in SINR, so interpolating a measured curve along SINR would compare a staircase with a ramp and manufacture a gap between the knots. The airtime share $B_c/K_c$ is kept in \emph{both} rates, since it is most of the rApp's benefit and cancels in the difference. Against the MAC the twin is right: $126.6\%$ predicted vs $125.4\%$ scheduled, a gap of $+1.2\pm5.7$ points, whose bootstrap $95\%$ CI $[-0.7, +4.1]$ straddles zero (Fig.~\ref{fig:realgap}, right). Against delivered goodput it is not: median $-6$ points, CI $[-22.2, -4.3]$, excluding zero---and that residual is the SINR-dependent emulator ceiling, not the abstraction. An OAI-grounded twin, corrected by one scalar $\kappa$, predicts what a real 5G MAC schedules, in absolute terms and in the ratios a controller consumes.

\begin{figure}[t]
\centering
\figorbox[0.72\columnwidth]{fig_real_fidelity_gap.pdf}{fig\_real\_fidelity\_gap.pdf}
\caption{Single gNB, matched bandwidth. Left: at matched operating SINR---the twin's link-level rate, that rate scaled by the TDD DL duty, OAI's scheduled MAC throughput, and delivered goodput. The twin-to-MAC ratio is constant ($\kappa=0.470\pm0.004$); the MAC-to-application spread is the emulator's ceiling and grows with SINR. Right: per-UE rApp steering gain. Against the MAC the twin sits on $y{=}x$ ($+1.2\pm5.7$ pts); against delivered goodput it does not.}
\label{fig:realgap}
\end{figure}

\subsection{Multi-cell closed loop, and what it does not yet measure}
\label{sec:multignb}
Two real cells (CU plus two F1 DUs, PCI~0/1) match a two-sector twin. The rApp's decision executes as a real inter-cell handover via scripted CU telnet, PDU session preserved: the loop closes. \emph{Measuring} the delivered gain is another matter. Our first attempt reported the twin's per-UE gain ($66\%$) overshooting the delivered gain ($9.0\%$) by $57\pm165$ points, and blamed the DUs' clean channel. That does not survive a null control: give both cells the \emph{same} channel and the true effect is exactly zero, so whatever the harness reports is bias. It reports $-21.0\pm1.4\%$. The cause is timing---a 6\,s TCP flow started 3\,s after each handover, inside both slow start and the link-adaptation ramp---and since \texttt{trigger\_f1\_ho} always moves DU0${\to}$DU1, measurement order and cell identity were perfectly confounded, so the reported sign followed whichever direction the rApp chose (on identical cells the probe yields mean DL MCS $16.5$ on PCI~0 vs $10.7$ on PCI~1). Probing at TCP's plateau gives $-2.4\pm2.9\%$, consistent with zero, MCS~28 on both. We withdraw the $57\pm165$ number.

With the instrument fixed we ground the channel (Bridge~3). Injection works: every layout's UE log confirms both per-cell models carrying the twin's ray-traced gains, and the UE reaches MCS~27 on the stronger cell. Two obstacles remain, neither the twin's. With channel models on both rfsim connections the loop destabilises---over three layouts two F1 handovers never complete, the third only through $70$ UE re-synchronisations. And rfsim's noise floor is a config knob, not $kTB\cdot\mathrm{NF}$: injection sets the \emph{relative} per-cell gain, not the absolute SINR, and sweeping the floor $-22$ to $-2$\,dB gives operating MCS $8,19,14,18,14,0$---non-monotone, hence not invertible. At the idle-neighbour condition a one-UE stack presents, the twin puts both cells at MCS~28 and predicts \emph{no} link-level difference in two of three layouts. A controlled multi-cell gap needs absolute operating-point matching, which this emulator cannot deliver.


\subsection{Coverage, outage, and load sensitivity}
Mean UE outage is $17.8\%$ at full neighbour load (Berlin $10.4\%$, Manhattan $17.1\%$, Paris $25.8\%$). What predicts a layout's rApp gain? Pooled over 24 runs, median gain correlates weakly with peak cell load ($r{=}0.29$) and negligibly with baseline outage ($r{=}0.05$). Both are unsafe---a correlation over three cities can be produced by the between-city contrast alone. Recomputed \emph{within} scene, neither survives: peak load gives $r{=}{+}0.46$, ${-}0.08$, ${+}0.18$ and outage ${-}0.42$, ${-}0.29$, ${+}0.08$ (Berlin, Paris, Manhattan), changing sign across scenes and, at eight seeds each, indistinguishable from zero. \textbf{We claim neither predictor.} An earlier version of this experiment---sub-band CFR, EESM applied band-wide, unseeded solver---reported a confident $r{=}{-}0.42$ on outage and dismissed load imbalance ($R^2{=}0.04$). The full band and a seeded solver dissolved it.

What does hold is the regime separation: lowering neighbour load from $1.0$ to $0$ raises mean throughput $1.68$--$1.95\times$, separating a coverage-limited tail steering cannot fix from an interference-limited body where load balancing acts.

%=============================================================================
\section{Discussion and Limitations}
\label{sec:disc}
%=============================================================================
We state these plainly. (1)~The loop is closed on one gNB (link calibration) and on two cells with scripted F1 handovers; autonomous E2-RC steering is future work. (2)~Bridge~2 awaits multi-cell rfsim validation and carries no fidelity claim. (3)~The single-gNB numbers are measured against \emph{this} emulator, which delivers only $24.9$ of the $88.9$\,Mbps its own MAC schedules at MCS~28; the twin's error is bounded by the MAC line, which is why we headline $\kappa$, not the $80\%$. (4)~Per-cell injection is verified and rfsim sums the DUs' waveforms---inter-cell interference is physically combined, superseding the noise-floor rise the analytical engine uses---but the multi-cell measurement does not close (Sec.~\ref{sec:multignb}), and the injected taps carry each link's gain, not its delay profile. (5)~The two constants are established for \emph{one} controller (CIO steering) on a proportional-fair objective; whether one scalar $\delta$ suffices for controllers reading MCS directly---link adaptation, MCS-aware schedulers---is open. (6)~Static full-buffer snapshot, single carrier, downlink; validation is stack-versus-model. Mobility, uplink, field calibration and real cell-site databases are next.

%=============================================================================
\section{Conclusion}
%=============================================================================
We asked how much of a real physical layer a 5G digital twin needs. For a proportional-fair traffic-steering rApp, two constants suffice. A \emph{fitted} implementation margin $\delta{=}4.71$\,dB cuts an uncalibrated twin's reporting error from $39.7\%$ to $1.6\%$ and its control regret from $0.767$ to $0.033$ nats---while a merely \emph{assumed} $2$\,dB margin does not ($0.732$), so it is $\delta$'s value, not a margin's presence, that matters. A derate $\kappa{=}0.470\pm0.004$---constant to $0.8\%$ across MCS 0--28---converts the twin's link rate into what a real 5G MAC schedules; being constant it cancels in ratios, so the twin's \emph{relative} steering-gain predictions are already right to $1.2\pm5.7$ points. Below the MAC lies an $80\%$ shortfall growing with SINR: the software radio's compute ceiling, not the twin's error. Two methods earn their place beside the constants. The customary ``decisions differ'' statistic ($25$--$38\%$ for every surrogate) cannot separate a faithful twin from an unfaithful one; only regret can. And a twin-versus-real comparison must be run against a null control---on two \emph{identical} real cells the natural goodput probe reports $-21.0\pm1.4\%$, which is how we came to retract a multi-cell gap we had published.

%=============================================================================
\section*{Reproducibility}
%=============================================================================
Everything is released under Apache-2.0. \textbf{Software:} OAI \texttt{develop}@\texttt{0dab3bb} (tag \texttt{2026.w27}); Sionna~RT $2.0.1$. $\Phi_{\mathrm{OAI}}$: sweep SNR\,$\times$\,MCS through \texttt{nr\_dlsim} (106 RB, 12 frames/point), taking at $0.2$\,dB granularity the lowest SNR whose first-transmission BLER meets the $10\%$ target, over $-2.6$ to $21.4$\,dB. Bridge~1 keeps a fixed count of the 4 strongest causal taps, unit-energy normalised, gain in \texttt{path\_loss\_dB}; Bridge~3 keeps one tap ($38.16$\,MHz resolves delay only to $26$\,ns). \textbf{Propagation:} specular reflection and refraction, interaction depth 3; diffuse scattering and edge diffraction disabled---the runtime--fidelity point that keeps a 24-run sweep on one CPU server, so we claim no diffuse-scattering fidelity. The solver is seeded from the layout seed: its path sampling is stochastic and an unseeded solve drifts every downstream number by ${\sim}1\%$. \textbf{Runtime:} a layout costs $13.1$\,s ray tracing $+\,4.7$\,s coordinate ascent; the sweep plus load sweep, $703$\,s, CPU-only. \textbf{Rebuilding:} \texttt{paper\_stats.py} recomputes every statistic here from committed CSV/JSON, \texttt{paper\_figures.py} every figure, and \texttt{paper/build.sh} fails if page count, overfull boxes or undefined references regress.

\balance
% IEEEtran's default inter-entry spacing; tighten it rather than cut a citation.
\def\IEEEbibitemsep{0pt plus .5pt}
\bibliographystyle{IEEEtran}
\bibliography{refs}

\end{document}
